\documentclass[11pt]{article}
\usepackage{amsmath,amssymb,mathtools}
\usepackage[margin=1in]{geometry}
\usepackage{xcolor}
\usepackage{booktabs,tabularx,longtable}
\usepackage{enumitem}
\usepackage[font=small,labelfont=bf]{caption}
\usepackage{microtype}
\usepackage{array}
\usepackage{upquote}
\usepackage[section]{placeins}
\usepackage[colorlinks=true,linkcolor=blue!50!black,citecolor=blue!50!black,urlcolor=blue!50!black]{hyperref}

\newcommand{\code}[1]{\texttt{#1}}
\newcolumntype{L}[1]{>{\raggedright\arraybackslash}p{#1}}
\newcolumntype{Y}{>{\raggedright\arraybackslash}X}
\newcommand{\sev}[1]{\textbf{#1}}
\setlist{leftmargin=1.4em,itemsep=1pt,topsep=2pt}
\renewcommand{\arraystretch}{1.12}

\title{The \code{@hss} Class: Audit and Plan for New Overloads\\[2pt]
\large what works, what is broken, and which MATLAB functions to add next}
\author{Arjun Sethi-Olowin (audit prepared with Claude)}
\date{October 2026}

\begin{document}
\maketitle

\section{Summary}
I read every file in \code{@hss/} and \code{+hssutil/}, ran the existing tests, and probed the class with
edge cases the tests do not cover. Everything was run on the repository code as it is (no file was modified),
in GNU Octave~8.4 through \code{octave\_compat}; items that depend on MATLAB itself are listed in
Section~\ref{sec:matlab} for you to confirm in R2025b.

\begin{itemize}
\item \textbf{The existing tests pass.} All 109 \code{test\_hss} methods (2160 checks; run in Octave through a small
stand-in for \code{matlab.unittest}) pass, except two checks that expect MATLAB's validator error identifier and
one test that needs \code{inter\_decompv2}'s \code{arguments} block. \code{mp\_test\_solver} reports 0 failures.
The solver, weighted min-norm, Tikhonov, products, transposes and extraction are correct on every input the
tests use.
\item \textbf{21 defects outside the tested inputs} (Table~\ref{tab:bugs}). Ten give a wrong answer with no
error or warning. The three to fix first:
\textbf{B06}, the randomized ID silently caps every off-diagonal rank at 120 (an exact HSS matrix with coupling
rank 70 is reproduced with relative error $0.79$);
\textbf{B02/B04}, \code{H(end,end)} returns \code{A(1,1)} (because \code{size(H)} is \code{[1 1]}) and
\code{H(mask,1:3)} with a two-entry mask returns a $64\times3$ matrix;
\textbf{B11}, a rank-deficient wide $H$ is ``solved'' with residual $4.1\times10^{-2}$ and no warning.
\item \textbf{Five design decisions} should come before new overloads (Section~\ref{sec:decisions}): the class
name clashes with hm-toolbox's \code{hss}; the property \code{size} blocks a \code{size} method; the recursive
object tree is slow to walk; the compression tolerance needs a definition; and whether to inherit from
\code{structuredmat}.
\item \textbf{New overloads} (Section~\ref{sec:overloads}): about 40 functions in four priority tiers, plus a list of
functions that should \emph{not} be overloaded and why. The six algorithms the plan depends on that are not
already in the repository (recompression, block extraction, a cell-array matvec, the adjoint of the stored
solve, log-determinant, exact Frobenius norm) were prototyped on the generator representation and checked
against dense references: 17 of 17 checks pass (Table~\ref{tab:proto}). The adjoint solve gives tall least
squares, \code{H'\textbackslash b} and \code{B/H} from the existing ULV factors, with no new factorization.
\item \textbf{Plan} (Section~\ref{sec:plan}): six phases, each with its files, tests and a check that ends it.
\end{itemize}

\noindent New files (nothing existing was changed), in \code{hss\_examples/audit\_2026\_10/}:
\code{hss\_audit\_probes.m} tests each defect (one line per bug; a line turns from \code{BUG} to \code{ok}
when the bug is fixed, so it doubles as a regression list; B21 runs in MATLAB only); \code{hss\_proto\_check.m} runs the prototype checks;
\code{proto/} holds the prototypes (\code{hssp\_*.m}).

% ---------------------------------------------------------------------------
\section{What the class does today}\label{sec:inventory}
\begin{table}[h]
\centering\small
\caption{Public entry points and the private code behind them, with the \code{test\_hss} tags that cover them.}
\label{tab:inventory}
\begin{tabularx}{\linewidth}{@{}L{2.9cm}YL{2.2cm}@{}}
\toprule
Entry point & What it does (implementation) & Test tags\\
\midrule
\code{hss(A, ...)} & Builds from a dense $A$: halves rows and columns in lockstep to a uniform depth; randomized
nested ID per off-diagonal block row/column (\code{inter\_decompv3}, sparse-sign sketch, threshold or fixed rank).
Options \code{blocksize, tol, k, decomp, cutrule, sizeA, powerits, \dots} & construction, compression\\
\code{H*X} & HSS times dense (\code{hss\_matvec}: up/down sweep, coefficients in a \code{dictionary}) & matvec, rectangular\\
\code{H1*H2} & HSS times HSS on the same tree; ranks add, no truncation (\code{hss\_matmat}) & matmat\\
\code{H\textbackslash B} & Square or wide: ULV min-norm solve (\code{hss\_ulvminnormsolve}); factors cached in
\code{H.factorcache}; tall: error & solve, rectangular, minnorm\\
\code{minnorm}, \code{tikhonov} & Weighted min-norm and Tikhonov via generator edits (\code{hss\_to\_gen},
\code{hss\_gen\_weight}, \code{hss\_gen\_augment}, \code{hss\_from\_gen}); own cache entries & minnorm\\
\code{H.'}, \code{H'} & Transpose / conjugate transpose by swapping generators & transpose\\
\code{H(I,J)} & Entry extraction (\code{subsref} $\to$ \code{extract}); \code{':'} supported & extract\\
\code{full}, \code{spy} & Dense matrix (\code{blockbuilder}); rank picture & full, construction\\
\code{clearfactors} & Frees the cached factors & --\\
\bottomrule
\end{tabularx}
\end{table}

\noindent Not overloaded today: \code{size}, \code{end}, \code{plus}, \code{minus}, \code{uminus}, scalar
\code{mtimes}, dense\code{*H}, \code{mrdivide}, \code{inv}, \code{det}, \code{norm}, \code{disp}, \code{subsasgn},
and every function in Section~\ref{sec:overloads}.

% ---------------------------------------------------------------------------
\section{Audit findings}\label{sec:bugs}
Severity: \sev{S1} wrong result without error or warning; \sev{S2} error or unusable result on valid input;
\sev{S3} performance or portability; \sev{S4} message or clean-up. Numbers come from Octave runs of
\code{hss\_audit\_probes.m} or of the experiments behind Section~\ref{sec:algorithms} (MATLAB timings will differ).

\begin{small}
\begin{longtable}{@{}L{0.75cm}L{0.6cm}L{5.4cm}L{7.6cm}@{}}
\caption{Defects, most severe first within each group. ``File:line'' points at the cause.}\label{tab:bugs}\\
\toprule
ID & Sev & Symptom (measured) & Cause and proposed fix\\
\midrule\endfirsthead
\toprule ID & Sev & Symptom & Cause and fix\\ \midrule\endhead
\midrule\multicolumn{4}{r}{\footnotesize continued}\\\endfoot
\bottomrule\endlastfoot
\multicolumn{4}{@{}l}{\emph{Construction}}\\
B06 & S1 & Exact HSS, coupling rank 70 (leaf block-row rank 140): rank 120 found, matvec rel.\ error
$0.79$; rank 100: $1.5$; rank 150: $2.8$. Coupling ranks 20 and 50 (block-row ranks 40, 100)
are exact. &
\code{inter\_decompv3.m:84,94}: threshold mode sketches with $\ell=\min(50,m,n)+10=60$ rows, so the sketch
can never see rank $>60$; ranks in $(60,120]$ are found only when the one-vector verification
(l.~380--406) triggers a redraw, at most twice ($\ell=60,120$). Fix:
adaptive range finder: if the detected rank reaches $\ell-p$ (the sketch is saturated), double $\ell$ and
resketch until it is not, or until $\ell\ge\min(m,n)$; then the verification is a safety net, not the rank
mechanism. Raise \code{hss:notLowRank} when $k=\min(m,n)$.\\
B08 & S1 & \code{hss(randn(64), cutrule=@(k) k-1)}: \code{full(H)} is $65\times65$. &
\code{hss\_constructor.m:162,241} trusts \code{cutrule}; the depth (l.~75--87) assumes halving, so uneven cuts
reach blocks of size 1 and then 0. Fix: require $1\le$\code{cutrule(k)}$\le k-1$, and compute the depth by
simulating \code{cutrule} itself, not halving.\\
B09 & S1 & \code{hss(randn(64), sizeA=[32 32])} is accepted; \code{H.size} is $[32~32]$. &
\code{sizeA} (l.~13) is an internal value exposed as an option. Fix: remove it from the options (or error when
it differs from \code{size(A)}).\\
B13 & S1 & A NaN entry: built without error; NaN spreads to most entries of every product (112 of 128 in
$H\mathbf 1$). &
No input check. Fix: \code{mustBeFinite} on \code{A}.\\
B14 & S2 & \code{hss(sparse A)} errors inside \code{inter\_decompv3} (Octave). In MATLAB the sketch $SA$ is
sparse, and sparse \code{qr} orders columns to reduce fill, not to reveal rank, so ranks would be wrong
(unverified). & Fix: \code{A = full(A)} block by block in the constructor (or accept sparse only through the
planned matvec-based constructor).\\
B15 & S2 & A matrix with fewer than $2\times$\code{blocksize} rows or columns (one leaf) has \code{Ir = Ic = []} and empty \code{blocksize};
\code{spy(H)} errors. & \code{hss\_constructor.m:95-97}. Fix: set \code{Ir=[1 m]}, \code{Ic=[1 n]},
\code{blocksize}, as \code{hss\_from\_gen} already does.\\
B20 & S3 & \code{+hssutil} has \code{sparsesign.c} and the compiled \code{sparsesign.mexmaca64} only:
\code{hss(A)} fails on Windows, Linux and Intel Macs unless someone runs \code{mex}. & Fix: add a plain
\code{+hssutil/sparsesign.m} (\code{octave\_compat/sparsesign\_octave.m} can serve) and use the MEX when
\code{exist(...,'file')==3}.\\
\multicolumn{4}{@{}l}{\emph{Indexing and size}}\\
B02 & S1 & \code{H(end,end)} returns \code{A(1,1)}; any \code{end} in an index is 1. &
No \code{end} method; MATLAB's default uses the object-array size, \code{[1 1]}. Fix: \code{end.m} returning
\code{H.size(k)} (needs B01).\\
B03 & S1 & \code{H(n+1,1)}, \code{H(0,1)}, \code{H(2.5,1)} return 0. &
\code{extract.m:33-38}: indices outside both halves match no \code{ismember} mask and keep the zero
initialization. Fix: validate indices in \code{subsref} (integer, $1\le i\le m$).\\
B05 & S1 & Internal nodes are not matrices: \code{H.A22(1:3,1:3)} returns all zeros; \code{H.A22*v} and
\code{H.A22\textbackslash v} error; \code{minnorm(H.A22,b,'Weight',w)} calls the tree non-uniform. &
Nodes keep global \code{Ir/Ic}, the global \code{levelcount}, and \code{isroot=false}. Fix: a
\code{subtree(H,path)} method that re-roots a copy (shift indices, set depth, add a cache); make the tree
properties \code{SetAccess=protected} and document \code{subtree} as the way to get a block (the probe
switches to \code{subtree(H,2)} once it exists).\\
B01 & S2 & \code{size(H)} is \code{[1 1]}; \code{numel}, \code{isscalar}, \code{length} follow. &
Property named \code{size}, no method. Fix: Decision D2.\\
B04 & S1 & \code{H(mask,1:3)} with a 64-entry mask holding 2 \code{true}s returns a $64\times3$ matrix (the mask
values 0/1 are used as indices); \code{H(k)} (linear) errors. & \code{extract.m:33-38}. Fix: convert logical
masks with \code{find};
map linear indices with \code{ind2sub}; support \code{H(:)} only if you want it (it densifies).\\
B19 & S3 & \code{H(1,end)} builds the whole $n/2\times n/2$ block: time $\times9.6$ from $n=2048$ to 8192
(probe). \code{H(I,J)}, $5\times4$ indices: $0.51$\,s at $n=8192$ vs $0.049$\,s for the prototype. &
\code{extract.m:45,50} call \code{blockbuilder}. Fix: \code{hssp\_extract} (Section~\ref{sec:extract}).\\
\multicolumn{4}{@{}l}{\emph{Solves}}\\
B10 & S4 & Square $H$ with a zero row and column: \code{H\textbackslash b} returns $x$ with residual
$4.8\times10^{-2}$ after only the generic ``matrix singular to machine precision'' warning of a triangular solve
(\code{hss\_ulvminnormsolve.m:271}); MATLAB's own \code{A\textbackslash b} behaves the same way. The user is not
told that $H$ is singular or where. & A singular $L_\tau$ is not checked (solver header). Fix:
after each \code{lq}, compare $\min_i|(L_\tau)_{ii}|$ with $n\,\epsilon\,\|H\|$ (estimate $\|H\|$ once by power
iteration); warn \code{hss:singular}; same with \code{rcond} for the root LU.\\
B11 & S1 & Wide $H$ with two equal rows, random $b$: residual $4.1\times10^{-2}$, no warning (open item in the
status doc). & \code{hss\_ulvminnormsolve.m:128} truncates the root SVD silently. Fix: warn
\code{hss:rankDeficient} when the numerical rank $r<m$, and report $r$.\\
B12 & S1 & \code{minnorm} with a singular leaf block in a cell \code{Weight} runs without error. &
\code{hss\_weightops.m:18-24} checks block sizes only. Fix: \code{rcond} check per block; reject NaN/Inf
vector weights.\\
B07 & S2 & \code{(H*H*H)\textbackslash b}: leaf rank 18 $>$ 16 leaf rows $\to$ \code{improperRanks}. &
\code{hss\_matmat} adds ranks with no truncation. Fix: \code{compress} (Section~\ref{sec:compress}); call it
inside \code{mtimes}/\code{plus} when a rank exceeds its block, or always with a tolerance option.\\
B16 & S2 & \code{H\textbackslash 3} for a one-row $H$ errors (``scalar mat multiplication''). &
\code{mldivide.m:12} treats any scalar \code{b} as scalar division. Fix: test \code{isscalar(b) \&\& H.size(1)
\textasciitilde= 1}.\\
\multicolumn{4}{@{}l}{\emph{Operators and messages}}\\
B17 & S2 & (a) \code{2*H}, \code{H*2}, \code{-H}, \code{H+H}, \code{x'*H} and (b) \code{B/H} all error. &
Not implemented. \code{x'*H} only needs \code{(H.'*x.').'}, since \code{transpose} exists. See
Section~\ref{sec:overloads}.\\
B18 & S4 & \code{H*v} with a wrong length: ``nonconformant arguments (op1 is 16x16, op2 is 21x1)''. &
No check in \code{mtimes}. Fix: check \code{size(X,1)==H.size(2)}, error \code{hss:mtimes:dimension}.\\
B21 & S2 & Two other \code{@hss} folders exist: hm-toolbox's (installed by \code{setup.m}) and
\code{Solver\_files/\allowbreak modifed\_HMtoolbox\_files2019/\allowbreak @hss}. Adding \code{structmats} ``with subfolders'' puts both on
the path, and \code{hss(A)} calls whichever comes first. & Fix: Decision D1.\\
\end{longtable}
\end{small}

\paragraph{Lower-priority notes.}
\begin{enumerate}[label=N\arabic*.]
\item \code{inter\_decompv3}'s header documents \code{ctype='rank'}, but the code tests for \code{'k'}; with
\code{'rank'} it silently uses rank 50. Validate \code{ctype}.
\item The ID verification uses one probe vector (\code{nverify = min(1,n)}, l.~307).
\item The constructor's \code{tol} is relative to each block's own largest pivot, so tiny blocks are resolved to
the same relative accuracy as large ones (more rank than a global accuracy needs). See D4.
\item \code{warning('Off-diagonal block may not be low-rank.')} has no identifier, so it cannot be silenced
selectively; it floods the test logs.
\item Debug output \code{disp('here rows')}, \code{disp('here cols')} remains (\code{hss\_constructor.m:326,350};
unreachable for trees the constructor builds).
\item Unused properties: \code{Q}, \code{S}, \code{minnormQ}, \code{nullorrange}; \code{lowrankrows/cols} are
meaningless after \code{mtimes} or \code{hss\_from\_gen}.
\item Children store the global \code{levelcount} (the leaf depth), not their own remaining depth (cause of part
of B05).
\item Tikhonov with $\lambda\gg\|H\|$: relative error of $x$ grows like $\epsilon\lambda/\|H\|$ ($6.7\times10^{-10}$ wide, $1.1\times10^{-9}$
tall at $\lambda=10^8$; $\le1.2\times10^{-13}$ for $\lambda\le10^2$). This is the conditioning of the augmented system,
not a bug; document it ($x\approx H^*b/\lambda^2$ there).
\item Lockstep row/column splitting: $64\times4096$ with \code{blocksize} 16 gives $16\times1024$ leaves; $300\times20000$ with
\code{blocksize} 200 gives no tree at all (one dense leaf). Warn when a leaf's aspect ratio is large; the real
fix is separate row and column depths (large change).
\item \code{hss\_matvec} needs \code{dictionary} (R2022b+). The cell-array version (\code{hssp\_matvec}) has no
such requirement; in Octave it is 20--115$\times$ faster and linear in $n$ (the Octave \code{dictionary}
stand-in is not; MATLAB numbers are for you to measure with \code{hss\_proto\_check}).
\item \code{test\_scalar\_multiply\_throws} and \code{test\_left\_multiply\_throws} assert current limitations; they
must change when the overloads land.
\end{enumerate}

% ---------------------------------------------------------------------------
\section{Decisions to make before adding overloads}\label{sec:decisions}
\begin{description}[leftmargin=1.2em,style=nextline]
\item[D1. Class name.] Three \code{@hss} folders can be on the path (B21). Options: (a) rename to
\code{hssmat} (matches \code{toeplitzmat}, \code{hankelmat}); (b) move it into a package
(\code{+structmats/\allowbreak @hss}, called as \code{structmats.hss}); (c) keep the name and drop hm-toolbox from the
path. I recommend (a): it is a one-time search-and-replace and removes the clash for every user.
\item[D2. The \code{size} property.] A \code{size} method is needed for \code{end}, \code{size(H,1)} and most
of MATLAB's generic code. A property and a method with the same name are at best confusing and, as far as I
know, not allowed in one class; check with a one-line \code{size.m}. Either way, rename the property (e.g.\
\code{sz}) in the same pass as D1. Do not overload \code{numel}: MATLAB's documentation advises classes with custom
indexing to leave it alone (it determines the number of outputs of indexing; \code{numArgumentsFromSubscript} is
the supported hook).
\item[D3. Storage: object tree or flat generators.] Every node is a full \code{hss} object with 25
properties (plus the hidden cache). \code{minnorm} and \code{tikhonov} already convert to generator cell arrays
(\code{hss\_to\_gen}) on every new factorization, and every prototype here works on those arrays.
Options: (a) keep the tree, add the new code on generators, convert on demand (least change); (b) store
the generators as the internal representation and build tree views only for \code{spy} and debugging (fastest:
the cell-array matvec is linear and much faster in Octave, B19 disappears, no \code{dictionary}). (b) touches
the constructor, solver and products, so plan it as its own phase (Phase~5) after the overloads exist and are
tested on (a).
\item[D4. Meaning of the tolerance.] For \code{compress}, \code{plus} and products, decide between a relative
tolerance per block (the constructor's current rule, N3) and an absolute one, $\tau\|H\|_2$ (hm-toolbox's
choice). The absolute rule bounds $\|H-\tilde H\|_2$ by a small, depth-dependent multiple of $\tau\|H\|_2$ and gives
smaller ranks; I recommend it, with $\|H\|_2$ from \code{normest}.
\item[D5. Inherit from \code{structuredmat}?] The repo's abstract class supplies \code{end} and naive fallbacks
that call \code{full}. Inheriting gives API consistency, but a fallback that densifies a $10^6\times10^6$
matrix after a warning is worse than an error. Recommendation: inherit only if every inherited fallback is
overridden with an HSS method or an explicit error.
\end{description}

% ---------------------------------------------------------------------------
\section{Proposed overloads}\label{sec:overloads}
Notation: $n$ the size, $r$ the maximum rank, $L=O(\log n)$ the depth. ``Proto'' marks algorithms checked in
\code{hss\_proto\_check.m}. Priority A is needed first, because other functions and generic MATLAB code rely on it.

\begin{small}
\begin{longtable}{@{}L{2.4cm}L{4.2cm}L{4.5cm}L{2.2cm}L{0.8cm}@{}}
\caption{Overloads to add, grouped by priority.}\label{tab:overloads}\\
\toprule
Function & MATLAB meaning for $H$ & HSS algorithm & Cost & Prio\\
\midrule\endfirsthead
\toprule Function & Meaning & Algorithm & Cost & Prio\\ \midrule\endhead
\midrule\multicolumn{5}{r}{\footnotesize continued}\\\endfoot
\bottomrule\endlastfoot
\multicolumn{5}{@{}l}{\emph{A. MATLAB semantics (indexing, size, display)}}\\
\code{size(H,dim)}, \code{length}, \code{ndims}, \code{isempty} & dimensions & read \code{sz} (D2) & $O(1)$ & A\\
\code{end} & last index in \code{H(i,j)} & \code{H.sz(k)} & $O(1)$ & A\\
\code{subsref} (rewrite) & \code{H(I,J)}: logical, linear, \code{end}, bounds checks & \code{hssp\_extract}:
partial bases restricted to $I,J$ & $O((|I|{+}|J|)rL$\newline$+|I||J|)$ & A, proto\\
\code{subsasgn} & \code{H(i,j)=v} & error \code{hss:readOnly} (an entry change is a rank-one change; offer
\code{lowrankupdate}) & -- & A\\
\code{disp}/\code{display} & show the matrix & one summary: size, depth, leaf sizes, min/mean/max rank, memory,
cached factors; never the whole tree & $O(n)$ & A\\
\code{isreal} & & all generators real & $O(nr)$ & A\\
\code{subtree} (new) & usable sub-block (B05) & re-rooted copy & $O(\text{size})$ & A\\
\multicolumn{5}{@{}l}{\emph{B. Arithmetic}}\\
\code{uminus}, \code{uplus} & $-H$, $+H$ & negate all $D_\tau$ and $B$ & $O(nr)$ & B, proto\\
\code{mtimes} & $sH$, $Hs$; $XH$ for dense $X$ & scale $D,B$; $XH=(H^{\mathsf T}X^{\mathsf T})^{\mathsf T}$ &
$O(nr)$; $O(nrp)$ & B, proto\\
\code{mrdivide} & $H/s$; $B/H$ & scale; $B/H=(H^*\backslash B^*)^*$: adjoint of $H$'s factors (wide/square
$H$), forward solve with $H^*$'s factors (tall $H$; min-norm, not MATLAB's basic solution) & $O(nr)$ & B\\
\code{plus}, \code{minus} & $H_1\pm H_2$ (same tree); $H\pm D$, $D$ diagonal or sparse diagonal; $H\pm s$ &
concatenate bases, $B=\mathrm{blkdiag}(B_1,B_2)$, then \code{compress}; diagonal: leaf $D_\tau$ only; scalar:
MATLAB adds $s$ to \emph{every} entry, a rank-one update (rank $+1$ everywhere). A mismatched tree is an error. &
$O(nr^2)$ & B, proto\\
\code{compress} (new) & re-truncate to tolerance, make proper & orthonormalize up, truncate down
(Section~\ref{sec:compress}) & $O(nr^2)$ & B, proto\\
\code{lowrankupdate} (new) & $H+UV^*$ & append $U,V$ to the bases at every level, then \code{compress} &
$O(n(r{+}p)^2)$\newline\hspace*{0pt} & B\\
\code{mpower} & $H^p$, integer $p$ & repeated squaring with \code{compress}; $p<0$: error until \code{inv} &
$O(nr^2$\newline$\cdot\log p)$ & B\\
\code{conj}, \code{real}, \code{imag} & & conj all generators; $\mathrm{real}(H)=(H+\bar H)/2$, then \code{compress} &
$O(nr^2)$ & B\\
\multicolumn{5}{@{}l}{\emph{C. Solves and scalar quantities}}\\
\code{mldivide}, tall $H$ & least squares $H^+b$ & factor $H^*$ (wide), apply the \emph{adjoint} of its min-norm
solve (Section~\ref{sec:adjoint}); cache under a second key & $O(nr^2)$, then $O(nr)$ & C, proto\\
\code{H'\textbackslash b} & $(H^*)^{-1}b$ or $(H^*)^+b$ & for square/wide $H$: adjoint of $H$'s own cached factors,
no new factorization. Needs \code{ctranspose}/\code{transpose} to hand $H$'s cache handle to the result with an
``adjoint'' flag, assigned \emph{after} the generator swaps (each set method otherwise gives $H'$ an empty cache) &
$O(nr)$ & C, proto\\
\code{logdet} (new), \code{det} & $\log|\det H|$, $\det H$ & sum of $\log|(L_\tau)_{ii}|$ over levels plus the root
LU; phase from $\det P_\tau\det Q_\tau$ (Section~\ref{sec:logdet}). \code{det} overflows for large $n$; document
\code{logdet} as the one to use & $O(nr^2)$ & C, proto\\
\code{norm(H,'fro')} & & exact from orthonormal generators & $O(nr^2)$ & C, proto\\
\code{norm(H)}, \code{normest} & $\|H\|_2$ & power iteration on $H^*H$ (matvecs) & $O(nrk)$ & C\\
\code{norm(H,1/Inf)} & & exact needs all entries ($O(n^2)$): use \code{normest1} with matvecs, say it is an estimate & $O(nrk)$ & C\\
\code{condest} & $\kappa_1$ estimate & \code{normest1} on $H$ and on $H^{-1}$ (cached solves) & $O(nrk)$ & C\\
\code{inv} & $H^{-1}$ & keep the error for now; later, HSS inversion from the ULV factors or construction from
solves (Phase 6). With cached factors, \code{H\textbackslash B} does what users need \code{inv} for & -- & C\\
\code{chol} & SPD factor & later: symmetric ULV (Xia); halves work, keeps symmetry & $O(nr^2)$ & C\\
\multicolumn{5}{@{}l}{\emph{D. Structure and spectra}}\\
\code{diag}, \code{trace} & & concatenate / sum leaf diagonals ($k=0$ only) & $O(n)$ & D\\
\code{tril}, \code{triu} & & zero the upper (lower) couplings, \code{tril} of each $D_\tau$; stays HSS &
$O(nr)$ & D\\
\code{blkdiag} & $\mathrm{diag}(H_1,H_2)$ & new root, zero couplings (equal depths) & $O(1)$ & D\\
\code{sum}, \code{mean} & & $H^{\mathsf T}\mathbf 1$ or $H\mathbf 1$ & $O(nr)$ & D\\
\code{eigs}, \code{svds} & a few eigen/singular values & wrap function handles: $Hx$; shift-invert with cached
solves & $O(nrk)$ & D\\
\code{issymmetric}, \code{ishermitian} & & randomized test $\|Hx-H^*x\|\le\tau\|H\|\|x\|$; say so in the help & $O(nr)$ & D\\
\end{longtable}
\end{small}

\paragraph{Not to overload, and why.}
\begin{itemize}
\item \code{numel}: changes the output count of indexing (D2).
\item \code{horzcat}, \code{vertcat}: $[H_1~H_2]$ is not HSS on any tree that respects both inputs (each block
row of $H_1$ couples to all of $H_2$ at full rank in general). Only \code{blkdiag} keeps the structure.
\item \code{eig}, \code{svd}, \code{lu}, \code{qr} in full: $O(n^3)$ or not structure-preserving; \code{eigs},
\code{svds} and the ULV factorization cover the uses.
\item \code{expm}, \code{sqrtm}, \code{logm}, \code{funm}: not cheaply HSS-preserving. A later option is a
rational approximation through shifted solves, each $O(nr^2)$.
\item Elementwise functions (\code{abs}, \code{exp}, \code{.\^{}}): destroy low rank. Hadamard product
\code{H1.*H2} is HSS with rank $r_1r_2$; add only if a use appears.
\item \code{double}, \code{single}, \code{sparse}: ambiguous for a compressed matrix; keep \code{full} as the one
explicit conversion.
\end{itemize}

% ---------------------------------------------------------------------------
\section{Algorithms behind the new functions}\label{sec:algorithms}
Generators as in \code{hss\_to\_gen}: leaf blocks $D_\tau$; leaf bases $U_\tau,V_\tau$; translations
$R_\tau,W_\tau$ with $\hat U_\tau=\mathrm{diag}(\hat U_{\alpha},\hat U_{\beta})R_\tau$ for children
$\alpha,\beta$ (same for $\hat V$); couplings $H(I_\alpha,J_\beta)=\hat U_\alpha B_{\alpha\beta}\hat V_\beta^*$.

\subsection{Recompression (\code{hssp\_orth}, \code{hssp\_compress})}\label{sec:compress}
\begin{enumerate}
\item \emph{Orthonormalize, bottom up.} At a leaf, $U_\tau=Q_\tau R_\tau$ (economic QR); replace $U_\tau$ by
$Q_\tau$ and push $R_\tau$ into everything that multiplies $\hat U_\tau$: the coupling where $\tau$ is the row
cluster ($B\leftarrow R_\tau B$) and $\tau$'s block of the parent translation. At a parent, QR the updated
translation and continue. Same for $V$ with $B\leftarrow BR_\tau^*$. Exact; ranks become
$\le$ block sizes, so this alone fixes \code{improperRanks}.
\item \emph{Truncate, top down.} With orthonormal bases, the off-diagonal block row of $\tau$ is
$\hat U_\tau S_\tau\,(\text{orthonormal})$, where $S_\tau=[\,B_{\tau,\text{sib}},\;R_\tau^{(\text{parent})}C_{p}\,]$
and $C_p=W_p^*S_p$ is the parent's truncated coefficient. Take the SVD $S_\tau=W\Sigma X^*$, keep $k'$ columns of
$W$, and update $\hat U_\tau\leftarrow \hat U_\tau W$ (leaf basis or translation columns),
$B_{\tau,\text{sib}}\leftarrow W^*B$, parent rows $\leftarrow W^*(\cdot)$. Then the same for columns.
\end{enumerate}
Prototype results: $H^3$ (leaf rank 18 $>$ 16 rows) compresses to rank 5 with relative error $1.8\times10^{-13}$,
and its solve agrees with dense to $2.1\times10^{-13}$; $H+H$ for an exact rank-10 HSS goes back to rank 10;
time 0.07, 0.15, 0.30, 0.59, 1.36\,s for $n=2^{12},\dots,2^{16}$ (Octave), i.e.\ linear.

\subsection{Block extraction (\code{hssp\_extract})}\label{sec:extract}
Sort the requested rows and columns into leaves; take the matching rows of each leaf basis; move up the tree with
$\hat U_p(I_p,:)=\mathrm{diag}(\hat U_\alpha(I_\alpha,:),\hat U_\beta(I_\beta,:))R_p$; at each node fill the two
coupling blocks $\hat U_\alpha(I_\alpha,:)B_{\alpha\beta}\hat V_\beta(J_\beta,:)^*$. Order and repeats are kept;
indices are validated (fixes B03). Times for \code{H(I,J)} with $|I|=5$, $|J|=4$: $0.020/0.131/0.512$\,s
(current) vs $0.008/0.023/0.049$\,s (prototype) at $n=1024/4096/8192$.

\subsection{Adjoint of the stored solve: tall least squares, \code{H'\textbackslash b}, \code{B/H}}\label{sec:adjoint}
For a wide or square $W$ of full row rank, the stored solve is a linear map $\mathcal S$ with $\mathcal S=W^+$
(memo~1). Each level of $\mathcal S$ is a linear map assembled from unitary maps ($P_\tau^*$, $Q_\tau^*$), the
isometries $\Omega_\tau^*$ (economic size reduction), triangular solves ($L_\tau^{-1}$), selections, one product
with the transformed matrix $\hat H$, and the next level. Taking the adjoint of this combination piece by piece
gives
$\mathcal S^*=(W^+)^*=(W^*)^+$, at the same $O(nr)$ cost, with $\hat H^*$ applied by the existing
\code{ctranspose}. Therefore:
\begin{itemize}
\item tall $H$ (full column rank): factor $W=H^*$ once; $H^+b=\mathcal S^*b$ (least squares);
\item square or wide $H$: $(H^*)^+b=\mathcal S^*b$ uses $H$'s own cached factors, no refactorization;
\item $B/H=\big((H^*)^+B^*\big)^*$ for wide or square $H$ (adjoint of $H$'s factors; this is exactly MATLAB's
\code{B/A=(A'\textbackslash B')'}). For tall $H$, $H^*$ is wide: use the \emph{forward} min-norm solve with $H^*$'s
factors, which gives $B H^+$; MATLAB would return a basic solution instead, the same deliberate deviation as
\code{H\textbackslash b} for wide $H$.
\end{itemize}
Prototype: tall $1024\times512$ least squares agrees with \code{pinv} to $6.5\times10^{-14}$ (complex,
Table~\ref{tab:proto}) and $7.4\times10^{-14}$ (real, separate run); square $H'\backslash b$ from $H$'s factors to $\le8\times10^{-13}$. A proof that the
adjoint equals $(W^*)^+$ is two lines from memo~1's Theorem (the solve is exactly $W^+$, and $(W^+)^*=(W^*)^+$);
it belongs in memo~1's open problems or as a short appendix.

\subsection{Log-determinant (\code{hssp\_logdet})}\label{sec:logdet}
For square leaves there is no size reduction, and one level gives $P^*HQ^*$ block lower triangular after the same
permutation of rows and columns (forced rows couple only to their own forced columns, through $L_\tau$).
Hence $\log|\det H|=\sum_\tau\log|\det L_\tau|+\log|\det H_{\text{red}}|$, recursively, ending with the root LU;
the phase is the product of $\det P_\tau\det Q_\tau/|\cdot|$ (the stored factors, $\hat H=P^*HQ^*$), the phases of
$\mathrm{diag}(L_\tau)$ and the root pivots. Prototype: absolute error $4.6\times10^{-13}$ to $1.1\times10^{-11}$
against dense LU on five matrices (real and complex random HSS, an exponential kernel with $\log|\det|=-2151$);
phase error $\le1.4\times10^{-12}$. Use: Gaussian-process likelihoods,
where $\log\det$ of a kernel matrix is the expensive part.

\subsection{Exact Frobenius norm, cell-array matvec}
After orthonormalization the blocks $D_\tau$ and $\hat U B\hat V^*$ are disjoint and the bases orthonormal, so
$\|H\|_F^2=\sum\|D_\tau\|_F^2+\sum\|B\|_F^2$ (error $5\times10^{-14}$). \code{hssp\_matvec} stores the up and
down coefficients in cell arrays indexed by level and node: $0.007$ vs $0.18$\,s at $n=4096$ and $0.10$ vs
$11.8$\,s at $n=65536$ in Octave (the second column is the Octave \code{dictionary} stand-in; measure MATLAB with
\code{hss\_proto\_check}).

\begin{table}[h]
\centering\small
\caption{Output of \code{hss\_proto\_check.m} (Octave 8.4). Errors are relative unless marked.}\label{tab:proto}
\begin{tabularx}{\linewidth}{@{}Yrr@{}}
\toprule
Check & Value & Limit\\
\midrule
$2H+(-3+i)H_B$ via \code{hssp\_scale}, \code{hssp\_plus} & $4.9\times10^{-18}$ & $10^{-13}$\\
\code{hssp\_orth}: change of the matrix & $2.6\times10^{-17}$ & $10^{-13}$\\
\code{hssp\_orth}: $\max_\tau\|U_\tau^*U_\tau-I\|$ (leaves) & $7.9\times10^{-16}$ & $10^{-13}$\\
\code{hssp\_compress}$(H^3)$: error & $1.8\times10^{-13}$ & $10^{-11}$\\
\code{hssp\_compress}$(H^3)$: leaf rank (was 18) & 5 & 8\\
\code{hssp\_compress}$(H^3)\backslash b$ vs dense & $2.1\times10^{-13}$ & $10^{-10}$\\
\code{hssp\_extract}: $H(I,J)$, any order, repeats & $7.5\times10^{-19}$ & $10^{-13}$\\
\code{hssp\_extract}: a full row & $7.8\times10^{-18}$ & $10^{-13}$\\
\code{hssp\_extract}: out-of-range index rejected & yes & yes\\
\code{hssp\_matvec} vs \code{H*x} & $1.9\times10^{-15}$ & $10^{-13}$\\
\code{hssp\_fro} vs sum of squares & $5.3\times10^{-14}$ & $10^{-13}$\\
tall least squares (adjoint) vs \code{pinv}, complex $1024\times512$ & $6.5\times10^{-14}$ & $10^{-11}$\\
square $H'\backslash b$ from $H$'s factors & $5.6\times10^{-13}$ & $10^{-11}$\\
\code{hssp\_logdet}: absolute error & $4.6\times10^{-13}$ & $10^{-9}$\\
\code{hssp\_logdet}: phase error & 0 & $10^{-12}$\\
\code{hssp\_logdet}, complex $H$: absolute error & $6.8\times10^{-13}$ & $10^{-9}$\\
\code{hssp\_logdet}, complex $H$: phase error & $1.4\times10^{-12}$ & $10^{-10}$\\
\bottomrule
\end{tabularx}
\end{table}

% ---------------------------------------------------------------------------
\section{Implementation plan}\label{sec:plan}
Each phase ends when its tests pass in MATLAB and the matching \code{hss\_audit\_probes} lines read \code{ok}.
New tests go into \code{test\_hss.m} under a tag per phase; small cases compare with dense, large cases
($n\ge2^{15}$) build from generators and never form the dense matrix.

\begin{description}[leftmargin=1.2em,style=nextline]
\item[Phase 0: decisions and safety (small).] D1--D5. Rename class/property, and update
\code{hss\_audit\_probes.m}, \code{hss\_proto\_check.m} and the tests in the same commit (they call \code{hss(...)}
and \code{H.size}). Fix B06 (adaptive ID), B08, B09, B13, B14, B15, B16, B18, B20, N1, N4, N5. Add
\code{mustBeFinite}, warning identifiers.
\emph{Done when:} probes B06, B08, B09, B13--B16, B18, B20, B21 read \code{ok}; the constructor tests still pass.
\item[Phase 1: MATLAB semantics (small--medium).] \code{size}, \code{end}, \code{length}, \code{ndims},
\code{isempty}, \code{isreal}, \code{disp}; rewrite \code{subsref} on \code{hssp\_extract} (logical, linear,
bounds); \code{subsasgn} error; \code{subtree}; \code{SetAccess=protected} on tree properties.
\emph{Done when:} B01--B05 and B19 read \code{ok} (B19: \code{H(1,end)} time ratio $\le6$ from $n=2048$ to 8192);
\code{H(end,end)}, \code{H(mask,:)}, \code{H(k)} match dense.
\item[Phase 2: arithmetic and compression (medium).] Port \code{hssp\_orth}/\code{hssp\_compress} as
\code{compress(H, tol)}; \code{uminus}, \code{uplus}, scalar and dense \code{mtimes}, \code{mrdivide} by a scalar,
\code{plus}/\code{minus} (HSS, diagonal, scalar), \code{lowrankupdate}, \code{mpower}, \code{conj},
\code{real}, \code{imag}; \code{mtimes(H1,H2)} compresses when a rank exceeds its block (B07).
\emph{Done when:} B07, B17a read \code{ok}; $\|\text{compress}(H)-H\|_2$ is a small multiple
of $\tau\|H\|_2$ on all test families (record the multiple in the test); ranks of $H+H$ equal those of $H$; compression time linear to $n=2^{17}$.
\item[Phase 3: solves and scalar quantities (medium).] Port \code{hssp\_ulv\_adjoint} into the solver file
(\code{solve\_level\_adjoint}); tall \code{H\textbackslash b}; \code{H'\textbackslash b} from cached factors (adjoint
flag on the transposed object's cache);
\code{B/H}; \code{logdet}, \code{det}; \code{norm} (fro, 2, estimates for 1/Inf), \code{normest},
\code{condest}; singularity and rank warnings (B10, B11), weight checks (B12).
\emph{Done when:} B10--B12 and B17b read \code{ok}; tall least squares within
$10\kappa\epsilon$ of dense QR on the memo-1 families; \code{logdet} within $10^{-10}|\log\det|$.
\item[Phase 4: structure and spectra (small).] \code{diag}, \code{trace}, \code{tril}, \code{triu},
\code{blkdiag}, \code{sum}, \code{mean}, \code{eigs}, \code{svds}, \code{issymmetric}, \code{ishermitian}.
\emph{Done when:} each matches dense at $n\le2048$; \code{eigs(H,k,'smallestabs')} reuses cached factors.
\item[Phase 5: flat generator storage (large, optional; D3).] Store generators; rewrite \code{mtimes},
the constructor output and the solver input on them; keep \code{hss\_from\_gen}/\code{hss\_to\_gen} for
compatibility. \emph{Done when:} the whole suite passes unchanged and \code{H*x} timing in MATLAB is no worse
than today at every $n$ (Octave suggests much better).
\item[Phase 6: later.] \code{inv} (HSS inversion), \code{chol} (symmetric ULV), a matvec/entry-based constructor
(\code{hss(@(x) A*x, @(x) A'*x, @(I,J) A(I,J), m, n)}, which never forms $A$), public
\code{hss.fromGenerators}, geometric splitting and rank-adaptive leaf sizes (open items in the status doc).
\end{description}

% ---------------------------------------------------------------------------
\section{To confirm in MATLAB}\label{sec:matlab}
\begin{enumerate}
\item Run \code{runtests('test\_hss')} (in \code{hss\_examples/tests}), then \code{hss\_audit\_probes} and
\code{hss\_proto\_check} (add \code{hss\_examples/audit\_2026\_10} to the path). Octave gives 20 \code{BUG} lines and
B20 \code{ok} (its copy has a \code{.m} fallback); in MATLAB expect 21--22 (B21 depends on your path) and 17/17
prototype checks.
\item B02: \code{H(end,end)} in MATLAB (expected identical to Octave: the default \code{end} sees a $1\times1$
object).
\item B14: what \code{hss(sparse A)} does in MATLAB.
\item B10: which warning MATLAB's triangular solve gives for the singular example (Octave: the generic
\code{Octave:singular-matrix}).
\item D2: whether MATLAB accepts a \code{size} method next to the \code{size} property.
\item Timing of \code{hssp\_matvec} vs \code{H*x} (printed by \code{hss\_proto\_check}); this decides how urgent
Phase~5 is.
\end{enumerate}

\end{document}
